% 
\RequirePackage{docswitch}
\setjournal{\flag}

\documentclass[\docopts]{\docclass}

% You could define the document class directly
%\documentclass[]{emulateapj}

% 
\usepackage{soul} 
\usepackage{amsmath}
\usepackage{xspace}
\usepackage{xifthen}
\usepackage[dvipsnames,svgnames]{xcolor} 

% General formatting
\newcommand{\ie}{i.e.\xspace}
\newcommand{\eg}{e.g.\xspace}
\newcommand{\etc}{etc.\xspace}
\newcommand{\etal}{et al.\xspace}
\newcommand{\vs}{vs.\xspace}
\newcommand{\super}[1]{\ensuremath{^{\textrm{#1}}}}
\newcommand{\sub}[1]{\ensuremath{_{\textrm{#1}}}}

\newcommand{\FIXME}[1]{{\bf \textcolor{red}{#1}}}
%\newcommand{\FIXME}[1]{{#1}}
\newcommand{\CHECK}[1]{{\bf \textcolor{orange}{#1}}}
%\newcommand{\CHECK}[1]{{#1}}
\newcommand{\COMMENT}[1]{{\it \textcolor{blue}{#1}}}
%\newcommand{\COMMENT}[1]{{#1}}
\newcommand{\NEW}[1]{{\textcolor{blue}{#1}}}
%\newcommand{\NEW}[1]{{#1}}

% Math
\mathchardef\mhyphen="2D
\newcommand{\vect}[1]{\boldsymbol{#1}}
\newcommand{\roughly}{\ensuremath{ {\sim}\,} }
\newcommand{\gtr}{\ensuremath{ {>}\,} }
\newcommand{\less}{\ensuremath{ {<}\,} }
\newlength{\dhatheight}
\newcommand{\doublehat}[1]{%
    \settoheight{\dhatheight}{\ensuremath{\hat{#1}}}%
    \addtolength{\dhatheight}{-0.35ex}%
    \hat{\vphantom{\rule{1pt}{\dhatheight}}%
    \smash{\hat{#1}}}}
\newcommand{\code}[1]{\texttt{#1}\xspace}
\newcommand{\dd}{\ensuremath{\rm d}}
\newcommand{\var}[1]{\ensuremath{#1}\xspace}


% Referencing 
\newcommand{\secref}[1]{Section~\ref{sec:#1}}
\newcommand{\appref}[1]{Appendix~\ref{app:#1}}
\newcommand{\tabref}[1]{Table~\ref{tab:#1}}
\newcommand{\tabrefs}[2]{Tables~\ref{tab:#1} and \ref{tab:#2}}
\newcommand{\figref}[1]{Figure~\ref{fig:#1}}
\newcommand{\figrefs}[2]{Figures~\ref{fig:#1} and \ref{fig:#2}}
\newcommand{\eqnref}[1]{Equation~\eqref{eqn:#1}}

% Astronomy
\newcommand{\LCDM}{\ensuremath{\rm \Lambda CDM}\xspace}
\newcommand{\ra}{{\ensuremath{\alpha_{2000}}}\xspace}
\newcommand{\dec}{{\ensuremath{\delta_{2000}}}\xspace}
\newcommand{\glon}{{\ensuremath{\ell}}\xspace}
\newcommand{\glat}{{\ensuremath{b}}\xspace}

% Units
\newcommand{\unit}[1]{\ensuremath{\mathrm{\,#1}}\xspace}
\newcommand{\Gyr}{\unit{Gyr}}
\newcommand{\MeV}{\unit{MeV}}
\newcommand{\GeV}{\unit{GeV}}
\newcommand{\TeV}{\unit{TeV}}
\newcommand{\degree}{\ensuremath{{}^{\circ}}\xspace}
\newcommand{\mas}{\unit{mas}}
\newcommand{\amin}{\unit{arcmin}}
\newcommand{\asec}{\unit{arcsec}}
\newcommand{\angstrom}{\unit{\AA}}
\newcommand{\um}{\unit{$\mu$m}}
\newcommand{\cm}{\unit{cm}}
\newcommand{\km}{\unit{km}}
\newcommand{\pc}{\unit{pc}}
\newcommand{\kpc}{\unit{kpc}}
\newcommand{\second}{\unit{s}}
\newcommand{\us}{\unit{$\mu$s}}
\newcommand{\photons}{\unit{ph}}
\newcommand{\photon}{\unit{ph}}
\newcommand{\sr}{\unit{sr}}
\newcommand{\Msolar}{\ensuremath{M_\odot}}
\newcommand{\Msun}{\ensuremath{M_\odot}}
\newcommand{\Mstar}{\ensuremath{M_{*}}}
\newcommand{\Lsolar}{\ensuremath{L_\odot}}
\newcommand{\Lsun}{\ensuremath{L_\odot}}
\newcommand{\Lstar}{\ensuremath{L_{*}}}
\newcommand{\Lum}{\ensuremath{ L }\xspace}
\newcommand{\cmcubes}{\ensuremath{\cm^{3}\second^{-1}}\xspace}
\newcommand{\magn}{\unit{mag}}
\newcommand{\mmag}{\unit{mmag}}
\providecommand{\deg}{}
\renewcommand{\deg}{\unit{deg}}
\newcommand{\kms}{{\km\second^{-1}}}
% 


\usepackage{graphicx}
\graphicspath{{./}{./figures/}}
\bibliographystyle{apj}


%%%%%%%%%%%%%%%%%%%%%%%%
%% Start the Document %%
%%%%%%%%%%%%%%%%%%%%%%%%

\begin{document} 

\title{Core Cosmology Library: Precision Cosmological Predictions for LSST}

\maketitlepre

\begin{abstract}
 
The Core Cosmology Library (CCL) provides routines which make precision theoretical predictions of cosmological observables for LSST. These routines have been validated to a $10^{-4}$ accuracy level against the results of the Code Comparsion Project. In the current version, predictions are provided for distances, angular auto- and cross-spectra of cosmic shear and galaxy number counts, and correlation functions of the same. The halo mass function is incorporated, and fiducial specifications for the expected LSST galaxy distributions and clustering bias are included. The facility is provided for a user-defined photo-z model. In this note, we explain the functionality of the first public release of the library. 

\end{abstract}

% Keywords for paper
%\dockeys{latex: templates, papers: awesome}

\maketitlepost

\section{Introduction}
\label{sec:intro}
In preparation for constraining cosmology with the Large Synoptic Survey Telescope (LSST), it is necessary to be able to produce theoretical predictions for the cosmological quantities which will be measured. The core cosmology library (CCL) aims to provide, in one library, predictions which are validated to a numerical accuracy of $10^{-4}$ or better for the purpose of constraining cosmology with LSST. By constructing a library with LSST in mind, it is possible to ensure that it is flexible, adaptable, appropriate, and validated for all cases of interest, as well as user-friendly and available for the needs of all working groups. 

The core cosmology library is written in C, which provides speed as well as allows the incorporation of the CLASS \cite{class} code. In the future, it is anticipated that python wrappers will be provided for improved ease of use. 

%This is a paper and note template for the LSST DESC \citep{Overview,ScienceBook,WhitePaper}.  
%Eventually it will be possible to switch between various \LaTeX\xspace styles for internal notes and peer reviewed journals templates.
%The base switch is between \code{aastex.cls} and \code{revtex.cls}; however, facilities are also provided for \code{emulateapj.cls} and \code{mnras.cls}.\footnote{The \code{mnras.cls} class file is a bit odd...}  
%Documents can be compiled using the provided \code{Makefile} with several options: \code{make apj}, \code{make apjl}, \code{make prd}, and \code{make mnras}. 
%There are some oddities when changing between templates, so please be patient while we try to work these out. 

%There are a number of useful \LaTeX\xspace commands predefined in \code{macros.tex}.
%Notice that the section labels are prefixed with \code{sec:} to allow the use of the \verb=\secref= command to reference a section (\ie, \secref{intro}).
%Figures can be referenced with the \verb=\figref= command, which assumes that the figure label is prefixed with \code{fig:}.
%In \figref{example} we show an example figure.
%You'll notice that the actual figure file is found in the \code{figures} directory. 
%However, because we have specified this directory in our \verb=\graphicspath= we do not need to explicitly specify the path to the image. 

%The \code{macros.tex} package also contains some conventional scientific units like \angstrom, \GeV, \Msun, etc. and some editorial tools for highlighting \FIXME{issues}, \CHECK{text to be checked}, \COMMENT{comments}, and \NEW{new additions}.
\section{Functionality}

\subsection{Permissible cosmological models}
\label{sec:cosmologies}
The user may be interested in producing cosmological constraints on $\Lambda$CDM parameters, or perhaps on some extension to the standard cosmological model. CCL handles this by allowing the user to create a cosmology object. (... more here ...) Functions exist to create general cosmologies, as well as for flat $\Lambda$CDM, curved $\Lambda$CDM, $\Lambda$CDM with and without neutrinos, wCDM, and the CPL model ($w_0$ + $w_a$).

%Similar to the figure before, here we have included a table of data from \code{tables/table.tex}.  
%Notice that again we are able to reference \tabref{example} with the \verb=\tabref= command using the \code{tab:} prefix. 
%Also notice that we haven't needed to specify the full path to the table because in the \code{Makefile} we include \code{./tables} directory in the \code{\$TEXINPUTS} environment variable.

%\begin{table}
  \begin{center}
  \caption{Example table. \label{tab:example}}
  %\begin{ruledtabular}
  \begin{tabular}{lccc}
\hline\hline
Column 1 & Column 2 & Column 3 &  Column 4 \\[3pt]  
     &    $\deg$     & $\kpc$   &  $\deg$ \\[4pt]
\hline
Obj1 & (0,0) & 10 & 0.1 \\
... & ... & ... & ... \\
ObjN & (0,0) & 10 & 0.1
\\\hline\hline
\end{tabular}
\end{center}
%\end{ruledtabular}
\end{table}

%\begin{\tabletype}{l ccccccc }
%\tablewidth{0pt}
%\tabletypesize{\tiny}
%\tablecaption{ An example table. \label{tab:example}}
%\tablehead{
%(1) & (2) & (3) & (4) & (5) & (6) & (7) & (8)\\
%Name & GLON,GLAT & Distance & $r_{1/2}$ & $\log_{10}(J_{\rm meas})$ & $\log_{10}(J_{\rm pred})$ & Sample & Refrence \\
% & (deg) & (kpc) & (pc) & $\log_{10}(\GeV^2 \cm^{-5})$ & $\log_{10}(\GeV^2 \cm^{-5})$ & & 
%}
%\startdata
%Bootes I                     & 358.08,69.62   & 66  & 189  & $18.8 \pm 0.2$ & 18.5           & I,S,C & ... \\
%\\
%...\\
%\\
%Willman 1                    & 158.58,56.78   & 38  & 19   & $19.1 \pm 0.3$ & 18.9           & I,S & ... \\
%\enddata
%{\footnotesize \tablecomments{ (1) The first column. (2) The second column ...}}
%\end{\tabletype}


\subsection{Distances}
\label{sec:distances}

The Hubble parameter is calculated via ... The comoving distance is calculated via ... (...more here...)

%If you are planning on committing your paper to github, it's a good idea to write your tex as one sentence per line. 
%This allows for an easier \code{diff} of changes.

\subsection{LSST Specifications}
\label{sec:specs}

CCL includes LSST specifications for the expected galaxy distributions of the full galaxy clustering sample and the lensing source galaxy sample. These enable the user to easily make predictions or forecasts for LSST.

The functional forms of the expected $\frac{dN}{dz}$ for clustering galaxies and lensing source galaxies are provided. Here, $\frac{dN}{dz}$ is the number density of galaxies as a function of spectroscopic redshift. 

In the case of lensing source galaxies, these forms are given in \cite{Chang2013}, wherein three different cases are considered: fiducial, optimistic, and conservative. All three are included in CCL, and are indicated via a label of DNDZ\_WL\_OPT, DNDZ\_WL\_FID, and DNDZ\_WL\_CONS as appropriate. The functional form of $\frac{dN}{dz}$ for lensing source galaxies is given as:
\begin{equation}
\frac{dN}{dz} \propto z^\alpha {\rm exp}\left(-\frac{z}{z_0}^\beta\right).
\label{dndz_src}
\end{equation}
The parameters, in the fiducial case, are given as $\alpha=1.24$, $\beta=1.01$, and $z_0=0.51$. In the optimistic case, this becomes $\alpha=1.23$, $\beta=1.05$, and $z_0=0.59$. The conservative case is given by $\alpha=1.28$, $\beta=0.97$, and $z_0=0.41$. 

For the case of the clustering galaxy sample, the functional form is given by \cite{ScienceBook}:
\begin{equation}
\frac{dN}{dz} \propto \frac{1}{2z_0}\left(\frac{z}{z_0}\right)^2 {\rm exp}\left(-\frac{z}{z_0}\right)
\label{dndz_clust}
\end{equation}
with $z_0=0.3$. The above $\frac{dN}{dz}$ for lensing sources in fact represents a subset of the $\frac{dN}{dz}$ for clustering.

In order to be incorporated into forecasts or predictions, the above expressions for $\frac{dN}{dz}$ must be normalized, and the value of $\frac{dN}{dz}$ must be provided in a given photometric redshift bin. Support is provided for the user to input a flexible photometric redshift model, as described further in Section \ref{sec:photoz} below. This takes the form of a function which returns the probability $p(z,z')$of measuring a particular photometric redshift $z$, given a spectroscopic redshift $z'$ and other relevant parameters. Also provided are functions to output $\sigma_z$ at a given redshift for both lensing sources and clustering galaxies, for the case in which the user wishes the assume a Gaussian photo-z model.

With this, $\frac{dN^i}{dz}$ of lensing or clustering galaxies in a particular photometric redshift bin $i$ is given by:
\begin{equation}
\frac{dN^i}{dz} = \frac{\frac{dN}{dz}\int_{z_i}^{z_{i+1}} dz' p(z,z')}{\int_{z_{\rm min}}^{z_{\rm max}}dz \frac{dN}{dz} \int_{z_i}^{z_{i+1}}dz' p(z, z')}
\label{photoz}
\end{equation}
where $z_{i}$ and $z_{i+1}$ are the photo-z edges of the bin in question. 

Finally, the expected (linear, scale-independent) bias of galaxies in the clustering sample is also provided. It is given by \cite{ScienceBook}:
\begin{equation}
b(z) = \frac{0.95}{D(z)}
\label{clustbias}
\end{equation}
where $D(z)$ in the linear growth rate of structure.

\subsection{Angular $C_\ell$'s}
\label{sec:cl}

In this section we will distinguish between {\sl observables} (inseparable quantities observed on the sky, such as number counts, shear or CMB temperature fluctuations) and {\sl contributions} to the total observed fluctuations of these observables (such as the main density term in number counts, redshift-space distortions, magnification, ISW, etc.).

\subsubsection{Exact expressions}
The angular power spectrum between two observables $a$ and $b$ can be written as:
\begin{equation}
 C^{ab}_\ell=4\pi\int_0^\infty \frac{dk}{k}\,\mathcal{P}_\Phi(k)\Delta^a_\ell(k)\Delta^b_\ell(k),
\end{equation}
where $\Delta^a$ and $\Delta^b$ are, using the terminology of CLASS, the transfer functions corresponding to these observables. Each transfer function will receive contributions from different terms. Currently CCL supports two observables, number counts and cosmic shear, with the following contributions:
\paragraph{\bf Number counts.} The transfer function for number counts can be decomposed into three terms: $\Delta^{\rm NC}=\Delta^{\rm D}+\Delta^{\rm RSD}+\Delta^{\rm M}$, where
\begin{itemize}
  \item $\Delta^{\rm D}$ is the standard density term proportional to the matter density:
        \begin{equation}
          \Delta^{\rm D}_\ell(k)=\int dz\,p_z(z)\,b(z)\,T_\delta(k,z)\,j_\ell(k\chi(z)),
        \end{equation}
        where $T_\delta$ is the matter transfer function. Note that CCL currently does not support non-linear or scale-independent bias. Here, $p_z(z)$ is the normalized distribution of sources in redshift (selection function). Thus CCL understand each individual redshift bin as a separate ``observable''.
  \item $\Delta^{\rm RSD}$ is the linear contribution from redshift-space distortions:
        \begin{equation}
          \Delta^{\rm RSD}_\ell(k)=\int dz\,p_z(z)\frac{(1+z) p_z(z)}{H(z)}T_\theta(k,z) j_\ell''(k\chi(z)),
        \end{equation}
        where $T_\theta(k,z)$ is the transfer function of $\theta$, the divergence of the comoving velocity field {\bf TODO: check that it's the comoving one - I'm 90\% sure of this}.
  \item $\Delta^{\rm M}$ is the contribution from magnification lensing:
        \begin{equation}
          \Delta_\ell^{\rm M}(k)=-\ell(\ell+1)\int \frac{dz}{H(z)} W^{\rm M}(z) T_{\phi+\psi}(k,z) j_\ell(k\chi(z)),
        \end{equation}
        where $T_{\phi+\psi}$ is the transfer function for the Newtonian-gauge scalar metric perturbations, and $W^{\rm M}$ is the magnification window function:
        \begin{equation}
           W^{\rm M}(z)\equiv\int_z^\infty dz' p_z(z')\frac{2-5s(z')}{2}\frac{\chi(z')-\chi(z)}{\chi(z')}.
        \end{equation}
        Here $s(z)$ is the magnification bias, given as the logarithmic derivative of the number of sources with magnitude limit.
          
        Note that CCL currently does not compute relativistic corrections to number counts {\bf TODO: add references}. Although these should be included in the future, their contribution to the total fluctuation is largely subdominant, and therefore it is safe to work without them for the time being {\bf TODO: add references?}.
\end{itemize}

\paragraph{\bf Cosmic shear.} The transfer function for cosmic shear is currently decomposed into two terms: $\Delta^{\rm SH}=\Delta^{\rm WL}+\Delta^{\rm IA}$, where
\begin{itemize}
  \item $\Delta^{\rm L}$ is the standard lensing contribution:
        \begin{equation}
          \Delta_\ell^{\rm L}(k)=-\frac{1}{2}\sqrt{\frac{(\ell+2)!}{(\ell-2)!}}\int \frac{dz}{H(z)} W^{\rm L}(z) T_{\phi+\psi}(k,z) j_\ell(k\chi(z)),
        \end{equation}
        where $W^{\rm L}$ is the lensing kernel, given by
        \begin{equation}
          W^L(z)\equiv\int_z^\infty dz' p_z(z')\frac{\chi(z')-\chi(z)}{\chi(z')}.
        \end{equation}
  \item $\Delta^{\rm IA}$ is the transfer function for intrinsic galaxy alignments. CCL currently supports the so-called ``linear alignment model'', according to which the galaxy inertia tensor is proportional the local tidal tensor.
        \begin{equation}
          \Delta_\ell^{\rm IA}(k)=\sqrt{\frac{(\ell+2)!}{(\ell-2)!}}\int dz\,p_z(z)\,b_{\rm IA}(z)\,f_{\rm red}(z)\,T_\delta(k,z)\,\frac{j_\ell(k\chi(z))}{(k\chi(z))^2}.
        \end{equation}
\end{itemize}

\subsubsection{The Limber approximation}
As shown above, computing each transfer function involves a radial projection (i.e. an integral over redshift or $\chi$), and thus computing full power spectrum consists of a triple integral for each $\ell$. This can be computationally intensive, but can be significantly simplified in certain regimes by using the Limber approximation, given by:
\begin{equation}
 j_\ell(x)\simeq\sqrt{\frac{\pi}{2\ell+1}}\,\delta\left(\ell+\frac{1}{2}-x\right).
\end{equation}
Thus for each $k$ and $\ell$ we can define a radial distance $\chi_\ell\equiv(\ell+1/2)/k$, and we will write the corresponding redshift as $z_\ell$. {\bf TODO: do we want to show why this is a good approximation with a plot?} This approximation works best for wide radial kernels and high multipoles.

Substituting this in the expressions above, it is possible to see that the The expressions above can be written as follows in the Limber approximation. First, the power spectrum
can be rewritten as
\begin{equation}
 C^{ab}_\ell=\frac{2}{2\ell+1}\int_0^\infty dk\,P_\delta\left(k,z_\ell\right)
 \tilde{\Delta}^a_\ell(k)\tilde{\Delta}^b_\ell(k).
\end{equation}
where
\begin{align}
 &\tilde{\Delta}_\ell^{\rm D}(k)=p_z(z_\ell)\,b(z_\ell)\,H(z_\ell)\\
 &\tilde{\Delta}_\ell^{\rm RSD}(k)=
 \frac{1+8\ell}{(2\ell+1)^2}\,p_z(z_\ell)\,f(z_\ell)\,H(z_\ell)-\\
 &\hspace{48pt}\frac{4}{2\ell+3}\sqrt{\frac{2\ell+1}{2\ell+3}}p_z(z_{\ell+1})\,f(z_{\ell+1})\,H(z_{\ell+1})\\
 &\tilde{\Delta}_\ell^{\rm M}(k)=3\Omega_{M,0}H_0^2\frac{\ell(\ell+1)}{k^2}\,
 \frac{(1+z_\ell)}{\chi_\ell}W^{\rm M}(z_\ell)\\
 &\tilde{\Delta}_\ell^{\rm L}(k)=\frac{3}{2}\Omega_{M,0}H_0^2\sqrt{\frac{(\ell+2)!}{(\ell-2)}}\frac{1}{k^2}\,
 \frac{1+z_\ell}{\chi_\ell}W^{\rm L}(z_\ell)\\
 &\tilde{\Delta}_\ell^{\rm IA}(k)=\sqrt{\frac{(\ell+2)!}{(\ell-2)!}}\frac{p_z(z_\ell)\,b_{\rm IA}(z_\ell)f_{\rm red}(z_\ell)H(z_\ell)}{(\ell+1/2)^2}
\end{align}

To be continued...

%And still not much.

\subsection{Correlation functions}
\label{sec:corr}

To produce correlation functions, we make use of the public code FFTlog (... cite ...). (... more here ... describe correlation function calculation).

%In conclusion, there are no conclusions.

%\begin{figure}
%\includegraphics[width=\columnwidth]{example.jpg}
%\caption{An example figure. \label{fig:example}}
%\end{figure}

\subsection{Halo mass function}
\label{sec:hmf}

The halo mass function is incorporated using the (Tinker 2008? cite) definition ... (...more here...).

\subsection{Photo-$z$ implementation}
\label{sec:photoz}
LSST galaxy redshifts will be obtained using photometry. However, analytic forms of galaxy redshift distributions are usually known in terms of spectroscopic redshifts, obtained as fits from a spectroscopic subsample (or extrapolation thereof). A model is therefore required for the probability of measuring a particular photometric redshift given a spectroscopic redshift. CCL allows the user to flexibly provide their own such photometric redshift model.

The user provides a function, which accepts a photometric redshift, a spectroscopic redshift, and a void pointer which points to a structure containing any further parameters of the photo-z model. The user then also provides this structure of parameters.

This user-provided model is incorporated when computing $\frac{dN}{dz}^i$ in photometric redshift bin $i$, as given by equation \ref{photoz}, above.


\section{Tests and validation}
\label{sec:tests}

Outputs of CCL have been validated against the results of the code comparison project down to a $10^{-4}$ or better accuracy level. (... more here ..., describe tests).

\section{Future functionality to be included}

In the future, we hope that CCL will include (... more here ...)

\section{Examples}
\label{sec:example}

%
This paper has undergone internal review in the LSST Dark Energy Science Collaboration. We thank the reviewers: Yao-Yuan Mao, Mariana Penna-Lima and Mike Jarvis for comments that halped improved this manuscript and the \ccl library overall. 

%DESC standard paper acknowledgements
%From https://github.com/LSSTDESC/desc-tex/blob/master/ack/standard.tex
The DESC acknowledges ongoing support from the Institut National de Physique Nucl\'eaire et de Physique des Particules in France; the Science \& Technology Facilities Council in the United Kingdom; and the Department of Energy, the National Science Foundation, and the LSST Corporation in the United States.  DESC uses resources of the IN2P3 Computing Center (CC-IN2P3--Lyon/Villeurbanne - France) funded by the Centre National de la Recherche Scientifique; the National Energy Research Scientific Computing Center, a DOE Office of Science User Facility supported by the Office of Science of the U.S.\ Department of Energy under Contract No.\ DE-AC02-05CH11231; STFC DiRAC HPC Facilities, funded by UK BIS National E-infrastructure capital grants; and the UK particle physics grid, supported by the GridPP Collaboration.  This work was performed in part under DOE Contract DE-AC02-76SF00515. 

%In addition
We would like to thank the organisers of the the DESC meetings at: CMU (July 2018), Oxford (July 2016), SLAC (Ferurary 2018, March 2016), and ANL (2015), and the LSST-DESC Hack Week organisers (CMU, November 2016), where this work was partly developed. We would also like to acknowledge the contribution of the participants of the Theory and Joint Probes Code Comparison Project, some of whom are among the \ccl contributors, for providing the benchmarks for testing CCL. We also acknowledge Louis Penafiel and Elizabeth Kimura, who developed the {\tt VARRIC} code to compare and visualize power spectra calculated by {\tt CCL} and {\tt CLASS}. Finally, we are grateful for the feedback received from other working groups of DESC, including Strong Lensing, Supernovae, Clusters and Photometric Redshifts.
We thank Peter Williams for making his ApJ bibstyle file available at \url{https://github.com/pkgw/tex-stuff/blob/master/yahapj.bst}. We are grateful to Katrin Heitmann and Earl Lawerence for discussions concerning the Cosmic Emulator. We are also thankful to the {\tt CLASS} authors and to Andrew Hamilton for making their codes available and allowing us to use them in this work. 
%

DA is supported by the Science and Technology Facilities Council (STFC) through an Ernest Rutherford Fellowship, grant reference ST/P004474/1. NEC acknowledges support from a Beecroft fellowship and a Royal Astronomical Society Research Fellowship. TT acknowledges funding from the European Union's Horizon 2020 research and innovation programme under the Marie Sk{l}odowska-Curie grant agreement No.\ 797794.MI acknowledges support by NSF under grant AST-1517768.


%{\it Facilities:} \facility{LSST}

\bibliography{main}

\end{document}
% 
